\clearpage
\section{Theoretical Background}
\label{sec:Background}
Supporting the choice between an LLM workflow and an agent requires three foundations: clear concepts, knowledge of what decision support must provide, and an overview of existing research. This chapter establishes these foundations. Section~\ref{subsec:theory/agentic-systems} places the choice within agentic systems. Section~\ref{subsec:theory/workflows-agents} defines the two orchestration paradigms, the task as the decision unit and four task dimensions that later structure the experiments. Section~\ref{subsec:theory/decision-support} derives the basis for the design objectives. Section~\ref{subsec:theory/related-work} reviews related work and positions the TADF within it.

\subsection{Agentic Systems and Automatisation}
\label{subsec:theory/agentic-systems}
Agentic systems pursue complex goals with minimal human intervention \parencite{acharyaAgenticAIAutonomous2025}. Their basic building block is a LLM that is augmented with retrieval, tools and memory \parencite{schluntzBuildingEffectiveAI2024}. With these capabilities, agentic systems plan steps, use tools to gather information and act on other systems, and keep context across steps \parencite{dwivediAgenticAISystems2026,sapkotaAIAgentsVs2026}. The related terms AI agent, agentic AI and agentic system are often used interchangeably \parencite{hughesAIAgentsAgentic2025}. Following \textcite{hughesAIAgentsAgentic2025}, this thesis uses the term agentic system for an interconnected structure that performs a business function. The term agent refers to one of the two orchestration paradigms defined in Section~\ref{subsec:theory/workflows-agents}.

In IS research, agency is a matter of degree. \textcite{holldackAgenticInformationSystems2026} place agentic information systems on a continuum. It ranges from assisting systems that act under significant human guidance to autonomous systems that manage complex tasks with minimal human intervention. Business process management describes a similar continuum for process execution. Process management systems have long executed predefined process models. Agentic systems can also execute such models, but they can additionally generate process variations, make contextual decisions and adapt workflows during execution. Agentic process management systems must therefore support processes that range from human-driven to fully autonomous \parencite{dumasAgenticBusinessProcess2026}.

Within an agentic system, control is also allocated between predefined code and the model. The orchestration layer of such a system decomposes high-level tasks into subtasks and coordinates them according to predefined or adaptive execution policies. In LangChain, chains provide a deterministic task structure, whereas agents add adaptive decision-making \parencite{sapkotaLangChainVsLangGraph2025}. Workflow and agent patterns can be combined within one system \parencite{schluntzBuildingEffectiveAI2024}. Deployed systems show such combinations. In 16 of 20 production case studies, practitioners defined the sequence of tasks upfront instead of relying on open-ended planning  \parencite{panMeasuringAgentsProduction2025}. An analysis of more than 6,000 workflows on the low-code platform n8n points in the same direction. Model calls are commonly embedded in broader structures with control logic, and a smaller share of workflows uses agent components with tools 
\parencite{tangCharacterizingLargeLanguage2026}.

This thesis examines the allocation of control for single tasks and therefore separates two levels for agentic systems. In the email process of Section~\ref{subsec:intro/motivation}, triage, context gathering, draft composition and draft verification are separate tasks. At the system level, these tasks are coordinated: results pass from one task to the next, and the system determines which task follows. At the task level, each task is executed by a component that combines model calls, program logic and tools. Within this component, control over the progression of the task can lie with predefined code or with the model. The TADF addresses this task-level choice. 

\subsection{LLM Workflow vs. Agent}
\label{subsec:theory/workflows-agents}

The terms workflow and agent are used with different meanings. OpenAI uses workflow for the sequence of steps that must be executed to meet a user's goal, whether by conventional software or by an agent \parencite{openaiPracticalGuideAgents2026}. A survey of agent workflows uses the term for procedures in which several agents make a sequence of decisions \parencite{yuSurveyAgentWorkflow2025}. Anthropic and LangChain instead distinguish workflows from agents by who controls the execution path. In a workflow, model calls and tools are orchestrated through predefined code paths. In an agent, the model directs its own process and tool use \parencite{schluntzBuildingEffectiveAI2024,Workflowsandgents}. OpenAI draws the same line: an application counts as an agent only if the model controls the execution of the workflow \parencite{openaiPracticalGuideAgents2026}.

This thesis adopts this control-based distinction. An LLM workflow coordinates model calls and tools through a path structure that is predefined in code. The structure can contain branches, loops and routing by a model. An agent leaves control over the progression of the task to the model. Within limits set by the surrounding program, the model selects the next action \parencite{schluntzBuildingEffectiveAI2024,Workflowsandgents}. It also decides when the task is complete. This thesis calls the two alternatives orchestration paradigms, or paradigms for short. The distinction concerns control over the execution path. It does not depend on the use of tools, the number of model calls or the implementation framework.

\subsubsection{Building Blocks and Architectural Differences}
\label{subsubsec:theory/control}

Both orchestration paradigms use the same building blocks: bounded model operations, program logic and tool operations. They differ in which block controls the progression of the task. In an LLM workflow, program logic fixes the order of the model operations, the conditions for branching or repeating a stage and the checks applied to results. Each model operation returns a defined output for a defined input, such as an extracted field or a draft text. \textcite{schluntzBuildingEffectiveAI2024} describe five workflow patterns: prompt chaining, routing, parallelization, orchestrator-workers and evaluator-optimizer. In each pattern, code specifies the stages and the permitted transitions, even where a model selects a route or plans subtasks within a stage. DSPy formalizes this composition as modules with declared input and output signatures, which ordinary control flow connects \parencite{khattabDSPyCompilingDeclarative2023}.

In an agent, the model controls the steps. In ReAct, the model alternates reasoning with actions, and each action returns an observation that becomes part of the context for the next step \parencite{yaoReActSynergizingReasoning2023}. The surrounding program provides the permitted tools, executes the selected actions, returns their results and enforces limits such as a maximum number of iterations \parencite{schluntzBuildingEffectiveAI2024}. Actions need not call external tools. The Cognitive Architectures for Language Agents (CoALA) framework also counts internal actions, such as reasoning over the working memory \parencite{sumersCognitiveArchitecturesLanguage2024}. Deciding whether to revise a draft or to submit it is therefore also a model-selected action.

Other design choices are independent of this distinction. The same model, tool interfaces and memory are available to both paradigms. CoALA distinguishes working memory for the current task from long-term episodic and procedural memory \parencite{sumersCognitiveArchitecturesLanguage2024}. What a system retains is therefore a separate design question. The same holds for the number of agents and for the agent style, such as ReAct prompting or native function calling. AgentArch varies both as separate configuration dimensions \parencite{bogavelliAgentArchComprehensiveBenchmark2025}. None of these choices determines who controls the progression of a task. Table~\ref{tab:theory-workflow-agent} summarizes the differences that follow from the allocation of control.

\begin{table}[htbp]
  \centering
  \caption{Control of task progression in LLM workflows and agents.}
  \label{tab:theory-workflow-agent}
  \small
  \begin{tabular}{@{}>{\raggedright\arraybackslash}p{\dimexpr0.20\linewidth-\tabcolsep\relax} >{\raggedright\arraybackslash}p{\dimexpr0.40\linewidth-1.5\tabcolsep\relax} >{\raggedright\arraybackslash}p{\dimexpr0.40\linewidth-1.5\tabcolsep\relax}@{}}
    \toprule
    \textbf{Aspect} & \textbf{LLM workflow} & \textbf{Agent} \\
    \midrule
    Task progression & Code specifies the stages and the transitions between them. & The model selects actions from the evolving task state. \\
    \addlinespace[2pt]
    Model contribution & Bounded outputs and choices within the predefined path. & Outputs and decisions about how to continue. \\
    \addlinespace[2pt]
    Steps and tool calls & Follow the conditions of the predefined path. & Depend on the model, up to the limits of the program. \\
    \addlinespace[2pt]
    Adaptation & Branches, repetitions and recovery follow coded conditions. & The sequence can change in response to observations. \\
    \addlinespace[2pt]
    Completion & The coded path defines when the task is complete. & The model can decide that the task is complete. \\
    \addlinespace[2pt]
    Traceability & Paths and transitions are known in advance; the trace shows which were taken. & Objectives and limits are known in advance; the sequence appears only in the trace. \\
    \bottomrule
  \end{tabular}
\end{table}

\subsubsection{Task as the Unit of Decision}
\label{subsubsec:theory/task}

Both orchestration paradigms apply to the same unit: the task. In Business Process Model and Notation (BPMN), a task is an atomic activity that is not broken down further at the chosen level of process detail \parencite{BPMN}. Building on this notion, this thesis defines a task as the smallest bounded unit of work with an identifiable input, a coherent goal, an assessable outcome and an progression over paths activities or tools. A task may require several model calls or tool operations. The definition does not assume that a task needs an LLM. Whether deterministic code suffices is part of the decision (Section~\ref{subsec:TADF/initial-artifact}).

Four further units structure the analysis. A task instance specifies a concrete input, the relevant initial state and the acceptance conditions. A task archetype groups tasks with recurring structural requirements. An implementation realizes one orchestration paradigm for a task archetype. An execution applies an implementation to a task instance.

\subsubsection{Task Dimensions Relevant to the Choice}
\label{subsubsec:theory/importance}
Which orchestration paradigm fits depends on the task. Practitioner guidance recommends LLM workflows for well-defined tasks and agents for open-ended problems (Section~\ref{subsec:intro/problem}). It also recommends the simplest solution that works, which may mean not building an agentic system at all \parencite{schluntzBuildingEffectiveAI2024}. Research on related design options supports task-dependent selection (Section~\ref{subsec:theory/related-work}). No single reasoning strategy performs best across tasks \parencite{zhouSelectthenSolveParadigmRouting2026}, and the effect of multi-agent coordination depends on the structure of the task \parencite{kimScienceScalingAgent2025}.

The literature points to four task dimensions that are candidates for explaining performance differences between the paradigms and task fit. This thesis later uses them to describe and structure tasks for the experiments (Section~\ref{subsubsec:TADF/taxonomy}). (1) Step predictability describes how far the required steps can be specified before execution. LLM workflows follow predefined code paths, whereas agents select their actions during execution \parencite{schluntzBuildingEffectiveAI2024,yaoReActSynergizingReasoning2023}. Research on workflow optimization likewise distinguishes structures fixed before deployment from structures that are selected, generated or revised for a particular run \parencite{yueStaticTemplatesDynamic2026}. (2) Information availability describes whether the required information is supplied with the task, held in known stores or obtained through interaction with the environment. CoALA separates retrieval from long-term memory from actions that ground the agent in external environments \parencite{sumersCognitiveArchitecturesLanguage2024}. FlowBench studies workflow knowledge that is supplied to an agent in advance \parencite{xiaoFlowBenchRevisitingBenchmarking2024}. (3) Output ambiguity describes how precisely an acceptable result can be specified. Evaluation metrics range from exact match to programs that balance several concerns \parencite{khattabDSPyCompilingDeclarative2023}. Iterative refinement by an evaluator is particularly effective when clear evaluation criteria exist and refinement provides measurable value \parencite{schluntzBuildingEffectiveAI2024}. (4) Error consequence describes the potential impact of an incorrect result in the intended application. In workplace settings, an agent error can send an email to the wrong person \parencite{stylesWorkBenchBenchmarkDataset2024}. Tasks that follow domain policies require reliable and consistent changes to the system state \parencite{yao$t$benchBenchmarkToolAgentUser2024}.

\subsection{Architecture Decision Support}
\label{subsec:theory/decision-support}
Choosing an orchestration paradigm for a task is a semi-structured decision. Evidence can show how well each paradigm suits a task. The weight of the competing qualities, however, depends on the deployment and on the judgment of the software architect (Section~\ref{subsec:intro/problem}). Decision support systems (DSS) target such decisions. A DSS has three components: data about the problem, models that evaluate alternatives, and a dialogue through which users enter information and inspect results \parencite{spragueFrameworkDevelopmentDecision1980}. The following paragraphs first state the aim of decision support. They then describe what decision support needs to reach this aim and to be used in practice.

Above all, decision support must make decisions measurably better. Improving decision quality is the goal of every DSS \parencite{kasperTheoryDecisionSupport1996}. For software engineering, \textcite{ruheSoftwareDesicionsupport} states this goal as a hypothesis: decision support enables more effective decisions of improved quality. The confidence of users does not show whether this goal is reached. According to \textcite{kasperTheoryDecisionSupport1996}, confidence in a decision often differs from its objective quality, and overconfidence is common. He also reports studies in which decision aids raised the confidence of their users but not the quality of their decisions. The value of decision support must therefore be shown by measuring the quality of the resulting decisions. This requires relevant metrics \parencite{article} and a reference point, such as decisions made without the support \parencite{ruheSoftwareDesicionsupport}.

To reach this aim, decision support starts from an explicit decision space. Techniques for architecture decisions describe a decision by its alternatives and by the quality attributes that these alternatives fulfill. Each alternative fulfills an attribute to a certain level, and this level is uncertain \parencite{falessiDecisionmakingTechniquesSoftware2011}. In the choice studied here, the alternatives are the two orchestration paradigms. A complete decision space can also include other found alternatives, which may suffice for some tasks (Section~\ref{subsubsec:theory/importance}). The quality attributes, such as task success, token cost and predictability, also provide the metrics for decision quality. Some attributes must reach a minimum level, whereas for others more is simply better \parencite{falessiDecisionmakingTechniquesSoftware2011}. Constraints go further and directly include or exclude parts of an architecture \parencite{amellerLinkingQualityAttributes2012}. Constraints and minimum levels therefore decide which alternatives are feasible, and preferences rank the feasible ones. 

Which feasible alternative fits best depends on the task. Task-technology fit is the correspondence between task requirements, individual abilities and the functionality of a technology. Only a technology that fits the task can improve individual performance \parencite{goodhueTaskTechnologyFitIndividual1995}. For the choice studied here, different task dimensions, deployment constraints and user requirements are assessed. The two paradigms offer different kinds of fit to these task requirements. The comparison therefore asks which paradigm provides the control that the task needs. Because the choice is made before implementation, it must rest on task characteristics that a software architect can state in advance.

This comparison must rest on evidence. \textcite{ruheSoftwareDesicionsupport} observes that decisions in software engineering often do not rely on empirically evaluated models and are not explained to those involved. Decision support in software engineering responds by combining models, measurement, empirical studies and simulation to evaluate and select alternatives \parencite{ruheSoftwareDesicionsupport}. Such evidence holds only for the tasks, technologies and conditions it covers. A recommendation must therefore name its evidence, separate observed findings from assumptions, and state the conditions under which it applies. 

The reasoning behind a recommendation must also be traceable. \textcite{ruheSoftwareDesicionsupport} accordingly includes an explanation component that makes solution alternatives transparent. \textcite{gregorExplanationsIntelligentSystems1999} classify explanations by their content, and two types matter here. A trace explains a recommendation by the data and rules used in the case. A justification links the rules to the knowledge from which they were derived. Experts use explanations more than novices to verify a recommendation or to resolve a disagreement, and justifications should increase trust and acceptance \parencite{gregorExplanationsIntelligentSystems1999}. Such verification is harder for advice from AI-based systems, which are less transparent than rule-based systems and whose errors are less predictable \parencite{jussupowAugmentingMedicalDiagnosis2021}. This thesis also considers this for reproducibility. Recording the inputs, the applied rules and their version ensures that the same inputs yield the same recommendation under the same version.

Decision support also extends to acting on a decision. \textcite{spragueFrameworkDevelopmentDecision1980} counts implementation among the phases of decision making that a DSS should support. For an orchestration paradigm, acting on a recommendation means building the LLM workflow or the agent. Recommendations therefore need guidance on how to implement the chosen paradigm. 

Finally, decision support helps only if professionals use it in their daily work. Fit with the task is not enough, because a technology improves individual performance only if it is also used \parencite{goodhueTaskTechnologyFitIndividual1995}. Users of a DSS can often decide for themselves whether to use it, so ease of use is particularly important \parencite{spragueFrameworkDevelopmentDecision1980}. According to \textcite{gregorExplanationsIntelligentSystems1999}, users read explanations only if the expected benefit outweighs the mental effort, and explanations requiring less effort are used more. They also point out that learning decisions and working constantly compete at work \parencite{lebovitzEngageNotEngage2022,gregorExplanationsIntelligentSystems1999}. Recommendations and implementation guidance must therefore be understandable for software architects and applicable with bounded effort. Together, these foundations describe what decision support for the choice of an orchestration paradigm must provide. Section~\ref{subsec:TADF/objectives} derives the design objectives of TADF from them.

\subsection{Related Work}
\label{subsec:theory/related-work}
Related research covers parts of the task-level choice. This section draws on the targeted review described in Section~\ref{subsec:intro/method}. It groups the studies into five streams. These streams differ in what they select, when they select it and on which evidence. Direct comparisons test both paradigms but offer no method for choosing between them. Studies of configurations select among agent designs only. Run-time methods select among implemented strategies, design guidance remains qualitative, and workplace benchmarks evaluate agents without a paired LLM workflow. No reviewed approach offers an evidence-based method for choosing between code control and model control for a single task before implementation.
 
The first stream is closest to the choice studied here. Two studies compare code control and model control directly. SourceCodeAgent follows the principle ``Blueprint First, Model Second''. An expert-defined procedure is written as code and run by a deterministic engine, while the LLM handles bounded subtasks but never decides the path \parencite{qiuBlueprintFirstModel2025}. Under the definitions of Section~\ref{subsec:theory/workflows-agents}, this is an LLM workflow. On the TravelPlanner benchmark with Claude Sonnet~4, it reached a final pass rate of 35.56\%, against 18.00\% for the multi-agent system ATLAS. The authors attribute the advantage to deterministic enforcement rather than a richer specification. On the explicit constraints of the user, however, the workflow performed comparably and produced more violations. They conclude that the payoff of code control scales with how procedurally well-defined and constraint-intensive a task is. For open-ended tasks, they consider general-purpose agents more suitable. \textcite{dennisInContextPromptingObsoletes2026} report the opposite for procedural dialogues. A complete procedure in the system prompt, followed by the model itself, reached higher judged quality than a LangGraph orchestrator with the same model. This held across three domains, but the prompted approach used more tokens. Both studies concern tasks with a defined procedure, yet they favor different paradigms. They differ in the kind of task and in what they measure. A defined procedure alone therefore does not settle the choice, and neither study offers a method for choosing a paradigm for a given task.
 
The second stream supports choosing the architecture for each task. It shows that the architecture of an agentic system affects its performance and that no single design fits every task and model. AgentArch evaluates 18 configurations on two enterprise use cases. They vary single-agent and multi-agent orchestration, ReAct and function calling, memory and thinking tools. The best configuration differs by model and task, which challenges one-size-fits-all designs \parencite{bogavelliAgentArchComprehensiveBenchmark2025}. \textcite{kimScienceScalingAgent2025} underline this and compare a single agent with four multi-agent architectures across 260 configurations, three model families and six benchmarks. Compared with the single agent, performance rose by up to 80.8\% on decomposable reasoning and fell by up to 70.0\% on sequential planning. The authors attribute the benefit of coordination to measurable task properties such as decomposability rather than to the number of agents. \textcite{gaoSingleagentMultiagentSystems2025} find that the accuracy advantage of multi-agent systems narrows as models improve and that simple tasks can suffer from overthinking. These studies compare configurations of agents, however, not code control with model control.
 
The third stream shows that selection pays off: choosing an execution strategy for each input can improve accuracy and reduce cost. \textcite{gaoSingleagentMultiagentSystems2025} route requests between single-agent and multi-agent systems. An LLM rater compares the difficulty of a request with a threshold. A cascade instead escalates a request only when the single-agent result fails a check. They report accuracy gains of 1.1 to 12\% at lower cost. Difficulty-Aware Agentic Orchestration (DAAO) starts from the observation that workflows fixed for a whole task over-process simple queries or underperform on complex ones. It estimates the difficulty of each query and composes a workflow from operators such as chain-of-thought, debate, review, ReAct and testing. DAAO reports up to 11.21\% higher accuracy at up to 36\% lower cost \parencite{suDifficultyAwareAgenticOrchestration2025}. Select-then-Solve compares six reasoning strategies across four models and ten benchmarks \parencite{zhouSelectthenSolveParadigmRouting2026}. A learned router reaches 53.1\% average accuracy, against 50.3\% for the best fixed strategy. Letting a model choose its own strategy helped only the strongest model, which still trailed the learned router. Select-Then-Decompose and LLM-as-Scheduler likewise choose or adapt strategies for each input, the latter with risk-tiered thresholds \parencite{liuSelectThenDecomposeEmpiricalAnalysis2025,xiangLLMasSchedulerAgenticWorkflow2026}. These approaches work during operation and choose among strategies that are already implemented, and learned routers need execution data for training. A software architect, in contrast, must decide before implementation which paradigm to build.
 
The fourth stream offers guidance for the same kind of design decision as this thesis. Besides the practitioner guidance discussed in Section~\ref{subsubsec:theory/importance}, the Agent Design Pattern Catalogue describes 18 architectural patterns with their context, forces and trade-offs. It proposes a decision model for selecting them and rests on a systematic literature review \parencite{liuAgentDesignPattern2025}. Its patterns concern how agents pursue goals and generate plans, not whether code or the model should control a task. In a survey of workflow optimization, \textcite{yueStaticTemplatesDynamic2026} recommend starting from a constrained static structure and adding flexibility only where the workload requires it. They reserve changes during execution for settings in which partial execution reveals information that a plan made in advance cannot anticipate. They call this recipe not universal. In their view, the field lacks a clear theory of when dynamic workflow generation is necessary and when static templates suffice. This guidance is qualitative and does not assess a given task against paired executions of both paradigms.
 
The fifth stream provides the realistic tasks and outcome checks on which evidence for the choice can be built. Benchmarks cover common business activities such as sending emails and scheduling meetings \parencite{stylesWorkBenchBenchmarkDataset2024}, realistic workflows of knowledge workers \parencite{boisvertWorkArenaCompositionalPlanning2025} and tasks in customer relationship management \parencite{huangCRMArenaProHolisticAssessment2025}. WorkBench counts a task as solved when the outcome of the actions matches a unique ground-truth outcome \parencite{stylesWorkBenchBenchmarkDataset2024}. $\tau$-bench likewise compares the final database state with an annotated goal state and adds dialogues with simulated users under domain policies. Even the strongest function-calling agent averaged a success rate below 50\% across its two domains and was inconsistent across repeated trials \parencite{yao$t$benchBenchmarkToolAgentUser2024}. The most competitive agent completed 30\% of the tasks in TheAgentCompany \parencite{xuTheAgentCompanyBenchmarkingLLM2025}. World of Workflows adds hidden workflows in enterprise systems whose effects cascade across connected databases \parencite{guptaWorldWorkflowsBenchmark2026}. On CRMArena-Pro, workflow execution was notably more tractable than the other business skills, with top agents exceeding 83\% success in single-turn tasks \parencite{huangCRMArenaProHolisticAssessment2025}. These benchmarks evaluate agents without a paired LLM workflow. FlowBench comes closest to a comparison with predefined control, but it varies only whether and in which format workflow knowledge is given to ReAct agents \parencite{xiaoFlowBenchRevisitingBenchmarking2024}. Section~\ref{subsubsec:TADF/taxonomy} uses these benchmarks as sources for the task archetypes of the experiments.